\documentclass[11pt,a4paper]{article}
\usepackage{operationsform}
\begin{document}
\OPStitle{記事構成・論理確認チェックリスト}{TPL-CHK-025}{0.2.0}
\OPSsection{1}{確認項目}
\OPScheck{1}{記事の目的、読者と中心となる問いが明確である}
\OPScheck{2}{見出し、導入、方法、結果、解釈、限界の各段落に説明上の役割がある}
\OPScheck{3}{重要な主張・数値・引用の箇所に出典IDを結び付けた}
\OPScheck{4}{文章の構成が時間順と因果の区別を壊していない}
\OPScheck{5}{異なる解釈や研究上の限界を必要な程度扱った}
\OPScheck{6}{各節の結論が根拠の確実性と適用範囲を超えない}
\OPScheck{7}{冒頭、見出し、結語と本文の確度が一致する}
\OPScheck{8}{編集上の未決事項と戻り先を記録した}
\OPSsection{2}{点検結果の記録}
\OPSblankshort{対象案件・対象版}
\OPSblankshort{担当者と点検日}
\OPSblanklong{未解決事項と対応}
\par\Needspace{4\baselineskip}\noindent\textbf{確認に用いた資料ID}\par\vspace{1mm}
{\fontsize{8.5}{13}\selectfont
\begin{tabularx}{\linewidth}{@{}XXXX@{}}
\textbf{対象・項目} & \textbf{内容・参照} & \textbf{担当・状態} & \textbf{備考}\\[2mm]
\strut & \strut & \strut & \strut\\[5mm]
\strut & \strut & \strut & \strut\\[5mm]
\end{tabularx}
}\vspace{2mm}\par
\end{document}
